\begin{proposition}[DP28: the direct higher-height intermediate estimate is false]\label{DP28}
For every odd $p$ and $h\ge3$, hypotheses (H1)--(H4) do not imply
\[
 v(\beta_{a,b})+\rho_h(a+b)\ge s-1.
\]
This refutes the stated substitution in the intermediate \path{LemmaS}, not the final derivative estimate.
\end{proposition}
\begin{proof}
Take $I(S,T)=1+S^pT$ and
\[
 \beta(X,Y)=\log I(E_h(X),E_h(Y)).
\]
Its coefficients lie in $\Q_p$. Since $E_h(\eta X)=\eta E_h(X)$ for $\eta^{q-1}=1$, the torus action cancels in $E_h(X)^pE_h(Y)$, proving (H2). Positive bidegrees and order at least $p+1$ give (H3). Formal logarithm and exponential identities give
\[
 H_X=\frac{I(F_h(S,S'),T)}{I(S,T)I(S',T)},\qquad
 H_Y=\frac{I(S,F_h(T,T'))}{I(S,T)I(S,T')}.
\]
All numerators are integral and all denominators are integral units, proving (H4).

Set $k=p^{h-1}<q$. In
\[
 \beta=\sum_{m\ge1}\frac{(-1)^{m-1}}m E_h(X)^{pm}E_h(Y)^m
\]
the coefficient of $Y^k$ can occur only for $m=k$, because $E_h(Y)=Y+O(Y^q)$. The coefficient of $X^{pk}=X^q$ in the remaining factor is one. Thus
\[
 \beta_{q,p^{h-1}}=\frac{(-1)^{p^{h-1}-1}}{p^{h-1}},\qquad
 v(\beta_{q,p^{h-1}})+\rho_h(q+p^{h-1})
 =2-h+\frac{p^{h-1}+1}{p^h-1}<s-1.
\]
The last inequality follows from $2-h\le-1$ and $(p-1)(p^{h-1}+1)<p^h-1$, equivalently $p<p^{h-1}$.
\end{proof}
